\documentclass[12pt]{book}
\usepackage[margin=0.8in]{geometry}
\parindent=0mm
\usepackage{tabularx,setspace, amssymb, amsmath, outlines, multicol, paralist, tikz, pgfplots, graphicx,enumitem, multienum, fixltx2e, fancyhdr,soul,lineno}
\pagestyle{fancy}
\everymath{\displaystyle}
\renewcommand{\indent}{\hspace{1cm}}
\newcommand{\longline}{\underline{\hspace{2in}} }
\newcommand{\shortline}{\underline{\hspace{1in}}}
%\renewcommand*\thesection{\arabic{section}}
%\renewcommand*\thesubsection{\thesection.\arabic{subsection}}
\renewcommand{\baselinestretch}{1.15}
\newcommand{\medium}{\textbf{Medium} }
\newcommand{\basic}{\textbf{Basic} }
\newcommand{\advanced}{\textbf{Advanced} }
\newcommand{\plane}{\begin{tikzpicture}[scale=0.6]
\draw[step=0.5] (0,0) grid (10,10);
\draw[ultra thick] (0,5) -- (10,5);
\draw[ultra thick] (5,0) -- (5,10);
\end{tikzpicture}}
\renewcommand*\tfrac[2]{
  \raisebox{-1em}{$\stackrel{\hbox{\underline{\smash{#1}}}}{\hbox{#2}}$}
}
\tikzset{every picture/.style={font issue=\footnotesize, scale=0.8},
	font issue/.style={execute at begin picture={#1\selectfont}},
	every path/.append style={line width=1 pt}
}

%\lfoot{\raisebox{-.25\height}{\includegraphics[scale=0.1]{logo}}{ \raisebox{0.25\height}{\small\copyright}\ Academic Student Center}}


\pgfplotsset{ every non boxed x axis/.append style={x axis line style=latex-latex}, every non boxed y axis/.append style={y axis line style=latex-latex},
axis x line=center, axis y line=center,
xmin=-5, xmax=5, ymin=-5, ymax=5,
axis line style = ultra thick, ticks=none,}

\begin{document}


\chapter{3 Strategies for Sentence Improvements}
\section{Warm-Up}

Are you having trouble remembering what types of errors are tested in the SAT Writing section?
Try this mnemonic device:

The most commonly tested and missed grammar points can be seen below. When you are answers sentence error or improvement questions, BE A CYCLOPS and always be keep one eye open for these most commonly missed grammar points. If you have already heard the BE A CYCLOPS mnemonic from another section, close your eyes and identify the grammar point that each letter refers to. Then, complete the exercise on the next page. 

\bigskip
\textbf{B is for “being”:} The word “being” is commonly heard in speech but does not usually make for the best sentences.

\bigskip
\textbf{E is for agrEEmEnt:} Identify the subject and the verb that is associated with the subject. The verb needs to match the subject in number and gender. This means that the subject and the main verb need to be both singular or both plural. 

\bigskip

\bigskip
\textbf{A is for awful verb tense:} Check when the action is happening and then if the given verb tense can be used to describe the time period that the action is happening. 

\bigskip

\bigskip
\textbf{C is for clause (aka commas towards the beginning of the sentence):} Clauses at the beginning of sentences have a description, then a comma, then more words. The description must be describing the first word after the comma. 

\bigskip
\textbf{Y is for you, me, and other pronouns:} If “you” is not in the underlined section, then it must be paired with “you” in the underlined section. If “one” or “someone” is not in the underlined section, then it must be paired with you in the underlined section. Also, make sure that pronouns like “it” or “they” clearly refer to the subject of the sentence. 

\bigskip
\textbf{C if for contrasts and other conjunction/connectors:} Words like “and” are used to add another idea, however, words like “but” are used to show differences between things. 

\bigskip
\textbf{L is for list:} If there is a list, all of the words must be the same part of speech and the same verb tense. 

\bigskip
\textbf{O is for “of” and commas that might separate the subject and the verb:} The verb ending is dependent on the singularity or plurality of the subject.

\bigskip
\textbf{P is for preposition:} Make sure the preposition matches the word before it. To combat this, learn your idioms!

\bigskip
\textbf{S is for short:} Is the sentence as short as it can be without changing the meaning? 

\vfill
\textit{Directions: Write 5 sentence error questions from any five different categories in the list above. At least one should have no error. Then, switch with someone in the class so that they can solve your questions.}

\begin{enumerate}
\item 
\vfill
\item 
\vfill
\item
\vfill
\item
\vfill
\item
\vfill
\end{enumerate}

\newpage


\section[Sentence Improvement]{Sentence Improvement Multiple Choice Questions}
Below is an example of a sentence improvement question:


\begin{enumerate}
\item The Boston Common is \ul{ older than it but still just as well-maintained as Central Park}.
\begin{enumerate}[label=(\Alph*)]

\item older than it but still just as well-maintained as Central Park
\item older than Central Park but just as well-maintained.
\item  older than Central Park; it is just as well-maintained.
\item older and it is just as well-maintained as Central Park.
\item just as comfortable as Central Park and it is older than it.

\end{enumerate}

\bigskip
Sentence Improvement Questions look extremely intimidating because the questions and the answer choices are relatively long. It looks like this would require a lot of reading; however, we are going to go through a strategy to decrease the amount of reading that you have to do but increase your accuracy and speed on the test. 

\section[Sentence Improvement]{How To Increase Your Score on Sentence Improvement Multiple Choice Questions}
\begin{enumerate}
\item Know the directions. The directions for the sentence improvement questions are as follows:

\textit{The following sentences test correctness and effectiveness of expression. Part of each sentence or the entire sentence is underline; beneath each sentence are five ways of phrasing the underlined material. Choice A repeats the original phrasing; the other four choices are different. Select the choice that completes the sentence most effectively. }

\textit{In making your selection, \textbf{follow the requirements of standard written English}; that is, pay attention to grammar, choice of words, sentence construction, and punctuation. \textbf{Your selection should result in the most effective sentence}-- clear and precise, without awkwardness or amiguity.}

\item Know the topics most commonly tested in sentence improvement questions.\\
Grammar: Subject-Verb Agreement, Pronoun Reference, Run-ons, Modifiers, Parallelism, Shift in Point of View, You and Me Errors,
Idioms, Hypothetical Statements, Punctuation
Phrasing: Redundancy, Conciseness, Sentence fragments

\textit{Note that these are slightly different from the sentence error questions and that a handful of topics, such as conciseness and punctuation, have been added.}

\item Read the original sentence. Determine if the original sentence has a blatant error or not. If so, circle the error in the question and think of 1-2 possible ways that you could correct this error.  

\item Move to the answer choices.  Eliminate the answer choices have the same error or that do not correct the error in the original sentence. (Usually ``A" and 1-2 other answer choices.)

\item Regardless of whether or not you were able to find an error in the original sentence, eliminate any other answer choices that are gramatically incorrect. 

\item This course of action will usually leave you with 1-2 answer choices that are grammatically correct. The directions for this section say ``Your selection should result in the most effective sentence", which means grammatically correct and also concise. When choosing between 2 gramatically correct sentences, the correct answer will be the more concise (shorter) answer choice that also  preserves the meaning of the original sentence. 

\end{enumerate} 

\textit{We will discuss these topics in more detail as the chapter progresses. }

\section{Punctuation}
Punctuation is more heavily emphasized in the sentence improvement and paragraph improvement sections than the sentence error sections. 
\begin{itemize}
\item When do you use a semi-colon on the SAT writing section? \hrulefill
\item How can you determine if a comma or semi-colon should be used? \hrulefill
\begin{enumerate}
\item Correct the following: Bob wanted to tip to the waiter, he could not do so because he had forgot his wallet at home.  
\end{enumerate}

\bigskip
\item What are the ways to correct a run-on sentence? \hrulefill

\bigskip
\item Correct the following original sentence in three different ways:
\begin{enumerate}
\item She hid quietly behind the sofa, if she was not discovered, then she would win the game of hide-and-seek. 

\bigskip
\item She hid quietly behind the sofa, if she was not discovered, then she would win the game of hide-and-seek. 

\bigskip
\item She hid quietly behind the sofa, if she was not discovered, then she would win the game of hide-and-seek. 
\end{enumerate}
\end{itemize}

\newpage



\section[Sentence Improvement]{Strategy 1: Identify the Error or Errors in the Original Sentence}
It is extremely advantageous if you can identify the error or errors in the original sentence before looking at the answer choices and then think of possible corrections. These can help you to 1) Eliminate answer choice ``A" as the best sentence and 2) eliminate incorrect answer choices with the same error as the original sentence quickly. 

\bigskip
\textit{Directions: Determine if the sentences below have an error in the underlined region and, if so, circle it. Write the type of error on the first line and a sample correction on the second line. If you don't that there is an error, write ``No error" as the error type and move to the next sentence. The first question has been done for you.}

\begin{enumerate}
\item The Boston Common is \ul{ older than it but still just as well-maintained as Central Park}.

\bigskip
\textbf{Type of error(s):} \ul{unclear pronoun, not concise} \\
\textbf{Sample correction:} \ul{older than but still just as well-maintained as Central Park.}

\bigskip
\item While most people detest high prices \ul{for food items, but organic food sells well despite the increased cost.}

\bigskip
\textbf{Type of error(s):} \hrulefill  \\
\textbf{Sample correction:} \hrulefill

\item With determination and dilligence, anyone can achieve a high score on the SAT test. 

\bigskip
\textbf{Type of error(s):} \hrulefill  \\
\textbf{Sample correction:} \hrulefill

\bigskip
\item The movie \ul{featured many well-respected actors and was winning many awards for} acting, directing, producing, and writing. 

\bigskip
\textbf{Type of error(s):} \hrulefill  \\
\textbf{Sample correction:} \hrulefill

\bigskip
\item \ul{Many educators believe that technology of the sort that} helps monitor student progress and deliever feedback to parents could be helpful in increasing test performance. 

\bigskip
\textbf{Type of error(s):} \hrulefill  \\
\textbf{Sample correction:} \hrulefill

\item After waiting an hour for her friend, the woman finally \ul{arrived in the theater donning a red dress.}

\bigskip
\textbf{Type of error(s):} \hrulefill  \\
\textbf{Sample correction:} \hrulefill

\item Overjoyed that he was accepted his first choice college, \ul{Stephen is currently being slightly ridiculous.}

\bigskip
\textbf{Type of error(s):} \hrulefill  \\
\textbf{Sample correction:} \hrulefill

\item \ul{Many people think that Americans take the right to vote for granted, and I think that it is the right of Americans to not exercise their right to vote.}

\bigskip
\textbf{Type of error(s):} \hrulefill  \\
\textbf{Sample correction:} \hrulefill

\item The bank robbers threatened the tellers by waving their guns, one of the criminals held a teller hostage until the police arrived. 

\bigskip
\textbf{Type of error(s):} \hrulefill  \\
\textbf{Sample correction:} \hrulefill

\item After a major political event such as September 11th, the president will address the nation, with his purpose being to inform and comfort the public.

\bigskip
\textbf{Type of error(s):} \hrulefill  \\
\textbf{Sample correction:} \hrulefill

\item Mary's secret, the whereabouts of the items that had been missing for weeks, \ul{were more compelling than Jeff's.}

\bigskip
\textbf{Type of error(s):} \hrulefill  \\
\textbf{Sample correction:} \hrulefill

\end{enumerate}

\newpage

\section{Strategy 2: Eliminating the Grammatically Incorrect Answer Choices}
On sentence improvements, the correct answer will be grammatically correct. Therefore, before
you see if an answer choice makes sense in the original sentence, determine if it is grammatically
correct. If not, then you can eliminate this right away.

\bigskip
\textit{Directions: Eliminate the answer choices that are grammatically incorrect for the following sentences. After you have elimintated an answer choice, write why it was incorrect on the line. The first one has been done for you.}

\begin{enumerate}
\item The Boston Common is \ul{ older than it but still just as well-maintained as Central Park}.
\begin{enumerate}[label=(\Alph*)]

\item older than it but still just as well-maintained as Central Park \textbf{ \ul{Eliminate because of ambiguous pronoun}}
\item older than Central Park but just as well-maintained. \longline
\item  older than Central Park; it is just as well-maintained. \textbf{ \ul{Eliminate because of ambiguous pronoun}}
\item older and it is just as well-maintained as Central Park. \longline
\item just as comfortable as Central Park and it is older than it. \textbf{ \ul{Eliminate because of ambiguous pronoun}}
\end{enumerate}

\bigskip
\textit{We now have two grammatically correct statements. Leave those alone-- we will choose the correct answer choice as part of the next exercise.}


\bigskip
\item While most people detest high prices \ul{for food items, but organic food sells well despite the increased cost.}

\bigskip
\begin{enumerate}[label=(\Alph*)]

\item        \hrulefill
\item   \hrulefill
\item    \hrulefill
\item    \hrulefill
\item   \hrulefill
\end{enumerate}

\bigskip
\item With determination and dilligence, anyone can achieve a high score on the SAT test. 

\bigskip
\begin{enumerate}[label=(\Alph*)]

\item        \hrulefill
\item   \hrulefill
\item    \hrulefill
\item    \hrulefill
\item   \hrulefill
\end{enumerate}

\bigskip
\item The movie \ul{featured many well-respected actors and was winning many awards for} acting, directing, producing, and writing. 

\bigskip
\begin{enumerate}[label=(\Alph*)]

\item        \hrulefill
\item   \hrulefill
\item    \hrulefill
\item    \hrulefill
\item   \hrulefill
\end{enumerate}

\bigskip
\item \ul{Many educators believe that technology of the sort that} helps monitor student progress and deliever feedback to parents could be helpful in increasing test performance. 

\bigskip
\begin{enumerate}[label=(\Alph*)]

\item        \hrulefill
\item   \hrulefill
\item    \hrulefill
\item    \hrulefill
\item   \hrulefill
\end{enumerate}

\bigskip
\item After waiting an hour for her friend, the woman finally \ul{arrived in the theater donning a red dress.}

\bigskip
\begin{enumerate}[label=(\Alph*)]

\item        \hrulefill
\item   \hrulefill
\item    \hrulefill
\item    \hrulefill
\item   \hrulefill
\end{enumerate}

\bigskip
\item Overjoyed that he was accepted his first choice college, \ul{Stephen is currently being slightly ridiculous.}

\bigskip
\begin{enumerate}[label=(\Alph*)]

\item        \hrulefill
\item   \hrulefill
\item    \hrulefill
\item    \hrulefill
\item   \hrulefill
\end{enumerate}

\bigskip
\item \ul{Many people think that Americans take the right to vote for granted, and I think that it is the right of Americans to not exercise their right to vote.}

\bigskip
\begin{enumerate}[label=(\Alph*)]

\item        \hrulefill
\item   \hrulefill
\item    \hrulefill
\item    \hrulefill
\item   \hrulefill
\end{enumerate}

\bigskip
\item The bank robbers threatened the tellers by waving their guns, one of the criminals held a teller hostage until the police arrived. 
\bigskip
\begin{enumerate}[label=(\Alph*)]

\item        \hrulefill
\item   \hrulefill
\item    \hrulefill
\item    \hrulefill
\item   \hrulefill
\end{enumerate}

\bigskip
\item After a major political event such as September 11th, the president will address the nation, with his purpose being to inform and comfort the public.

\bigskip
\begin{enumerate}[label=(\Alph*)]

\item        \hrulefill
\item   \hrulefill
\item    \hrulefill
\item    \hrulefill
\item   \hrulefill
\end{enumerate}

\bigskip
\item Mary's secret, the whereabouts of the items that had been missing for weeks, \ul{were more compelling than Jeff's.}

\bigskip
\begin{enumerate}[label=(\Alph*)]

\item        \hrulefill
\item   \hrulefill
\item    \hrulefill
\item    \hrulefill
\item   \hrulefill
\end{enumerate}


\end{enumerate}

\newpage
\section{Strategy 3: Picking the Most Clear and Concise Sentence}
In SAT sentence improvement problems, you should try to eliminate the 2-3 answer choices that are grammatically incorrect so that you are left with 2-3 other answer choices. The SAT sentence improvement section is looking for the "best" sentence, one that is concise and precise. In SAT world, this translates to the sentence that is not only grammatically correct AND concise. How does the SAT measure "conciseness"? By length. \\

Therefore, the answer choice that you are looking for is grammatically correct and short without changing the meaning of the original sentence. The latter part means that it can not be so short that it is missing a key part of the original sentence,
but this is not usually an issue on sentence improvement problems. \\

\textit{Directions: Go back to the previous exercise and look at the answer choices that you haven't yet eliminated. Find the shortest answer. Then, check if it preserves the original meaning of the sentence by reading this answer choice in place of the underlined part of the original sentence. If so, this answer choice is the correct answer, so you should mark it. The first sentence is done as an example.}

\begin{enumerate}
\item The Boston Common is \ul{ older than it but still just as well-maintained as Central Park}.
\begin{enumerate}[label=(\Alph*)]

\item older than it but still just as well-maintained as Central Park  \ul{Eliminate because of ambiguous pronoun}
\item older than Central Park but just as well-maintained. \textbf{\ul{Correct. It is grammatically correct and the most concise. }}
\item  older than Central Park; it is just as well-maintained.    \ul{Eliminate because of ambiguous pronoun}
\item older and it is just as well-maintained as Central Park. \textbf{\ul{Grammatically correct but not as concise as (B).}}
\item just as comfortable as Central Park and it is older than it. \ul{Eliminate because of ambiguous pronoun}

\end{enumerate}

\end{enumerate}

\newpage

\section[Sentence Improvement]{Practice with Sentence Improvement Multiple Choice Questions}
\textit{Directions: The following sentences test correctness and effectiveness of expression. Part of each sentence or the entire sentence is underline; beneath each sentence are five ways of phrasing the underlined material. Choice A repeats the original phrasing; the other four choices are different. Select the choice that completes the sentence most effectively. }

\textit{In making your selection, follow the requirements of standard written English; that is, pay attention to grammar, choice of words, sentence construction, and punctuation. Your selection should result in the most effective sentence-- clear and precise, without awkwardness or amiguity. \textbf{After you have selected the correct answer choice, you need to mark why the other four answer choices are incorrect, just like in strategies 1-3.}}

\begin{enumerate}
\item When asked why she became a journalist, the woman responded that she wanted to help tell people’s stories and \ul{ loved to be on camera.}

\bigskip
\begin{enumerate}[label=(\Alph*)]
\item    \hrulefill
\item    \hrulefill
\item   \hrulefill
\item   \hrulefill
\item   \hrulefill
\end{enumerate}

\end{enumerate}

\end{enumerate}

\end{document}